% !PRINTED 68
\weberpage{68}
%<<CONT
determine, we take into view the dependence of the function $t$ upon
$\omega_1, \omega_2$ as well, and, regarding $g_1, g_2$ as given numbers
independent of $\omega_1, \omega_2$, we denote it by
$t(u, \omega_1, \omega_2)$. It then appears that, if $h$ is an arbitrary
factor,
\[
t(h\,u,\ h\,\omega_1,\ h\,\omega_2)
\]
likewise satisfies the conditions (1), and that accordingly

\weq{2}{\frac{t(hu,\ h\omega_1,\ h\omega_2)}{t(u,\ \omega_1,\ \omega_2)} = f(\omega_1,\ \omega_2)}

is independent of $u$.

But the function $f(\omega_1, \omega_2)$ can always be put into the
form:
\[
f(\omega_1,\ \omega_2) = \frac{\varphi(h\omega_1,\ h\omega_2)}{\varphi(\omega_1,\ \omega_2)};
\]
one has only to put, if $t(u)$ does not vanish for $u = 0$:
\[
\varphi(\omega_1,\ \omega_2) = t(0,\ \omega_1,\ \omega_2),
\]
and, if $t(0) = 0$, say
\[
\varphi(\omega_1,\ \omega_2) = \omega_1\, t'(0,\ \omega_1,\ \omega_2)
\]
where $t'$ denotes the derivative of $t$ with respect to $u$. If, then,
one now denotes

\weq{3}{\frac{t(u,\ \omega_1,\ \omega_2)}{\varphi(\omega_1,\ \omega_2)}}

again by $t(u, \omega_1, \omega_2)$, one may add to the conditions (1)
the further one that for an arbitrary $h$

\weq{4}{t(u,\ \omega_1,\ \omega_2) = t(hu,\ h\omega_1\ h\omega_2).} % [sic] -- no comma between $h\omega_1$ and $h\omega_2$, as printed in the original

The function $t$ thus depends now upon only two variables, namely upon
the ratios $u : \omega_1 : \omega_2$. In (4) we put
\[
h = \frac{1}{\omega_1},\quad \frac{u}{\omega_1} = v,\quad \frac{\omega_2}{\omega_1} = \omega,
\]
and obtain

\weq{5}{t(u,\ \omega_1,\ \omega_2) = t\left(\frac{u}{\omega_1},\ 1,\ \frac{\omega_2}{\omega_1}\right) = \vartheta(v, \omega),}

wherein $\vartheta$ is a new function symbol.

In this way, then, a new function $\vartheta(v)$ is defined, of only
two variables $v, \omega$, the first of which is unrestrictedly
variable, while $\omega$, by the assumption made in \S\,16, has a
\gesperrt{positive imaginary} part; $v$ is called the \gesperrt{argument},
$\omega$ the \gesperrt{modulus} of the $\vartheta$-function.
